\section*{Linear circuits and sparse matrix products (original Section 1.1)}
Linear circuits are algebraic circuits which are only allowed to use addition and multiplication by a scalar. It is often convenient to consider graphs with labels on the edges as well: every internal node is an addition gate, and for $c\in\F$, an edge labeled $c$ from a vertex $v$ to a vertex $u$ denotes that the output of $v$ is multiplied by $c$ when feeding into $u$. Thus, every node computes a linear combination of its inputs. The size of the circuit is defined to be the number of edges in it. The depth of a circuit is the length (in edges) of a longest path from an input to an output.
A depth 2 linear circuit computing a linear transformation $A\in\F^{n\times n}$ consists of a bottom layer of $n$ input gates, a middle layer of $m$ gates, and a top layer of $n$ output gates. We assume, without loss of generality, that the circuit is layered, in the sense that every edge goes either from the bottom to the middle layer, or from the middle to the top layer. Indeed, every edge going directly from the bottom to the top layer can be replaced by a path of length 2; this transformation increases the size of the circuit by at most a factor of 2.
By letting $C\in\F^{m\times n}$ be the adjacency matrix of the (labeled) subgraph between the bottom and the middle layer, and $B\in\F^{n\times m}$ be the adjacency matrix as the subgraph between the middle and the top layer, it is clear that $A=BC$. Thus, a decomposition of $A$ into the product of two sparse matrices is equivalent to saying that $A$ has a small depth-2 linear circuit. This argument can be generalized, in exactly the same way, to depth-$d$ circuits and decompositions of the form $A=A_1\cdots A_d$, for constant $d$.
